\documentclass[dvipdfmx]{jsarticle}
\usepackage{graphicx}
\usepackage{url}

\title{ソフトウェア設計及び実験\\第10回（バージョン管理について）}
\author{学生番号 \ 6125010806　氏名 \ 清水 \ 統貴}
\date{2026年6月18日提出}

\begin{document}
\maketitle

\section{課題の目的}
%
% 1. 目的
% 2. 原理・理論
% 3. 方法
% 4. 結果
% 5. 考察
% 6. 参考文献
% 7. 付録
本課題の目的は以下のとおりである．\\
• IDEの使い方について理解する\\

% \section{原理}
% %

\section{課題1}
%
以下, 実際にsvnをつかった時の実行したコマンドと結果.\\
\begin{verbatim}
    c612501080@clpc019:~/soft_experiment/10$ mkdir repos
    c612501080@clpc019:~/soft_experiment/10$ svnadmin create repos
    c612501080@clpc019:~/soft_experiment/10$ mkdir work
    c612501080@clpc019:~/soft_experiment/10$ cd work
    c612501080@clpc019:~/soft_experiment/10/work$ mkdir myproject
    c612501080@clpc019:~/soft_experiment/10/work$ cd myproject
    c612501080@clpc019:~/soft_experiment/10/work/myproject$ echo This is readme file. > readme.txt
    c612501080@clpc019:~/soft_experiment/10/work/myproject$ cat readme.txt 
    This is readme file.
    c612501080@clpc019:~/soft_experiment/10/work/myproject$ svn import file:///usrhome/OTF25/c612501080/soft_experiment/10/repos/myproject -m "Initial import."
    追加しています              
    readme.txt
    リビジョン 1 をコミットしました。
    (元あったmyprojectを消去後)
    c612501080@clpc019:~/soft_experiment/10/work$ svn checkout file:///usrhome/OTF25/c612501080/soft_experiment/10/repos/myproject myproject
    A  myproject\readme.txt
    リビジョン 1 をチェックアウトしました。
    c612501080@clpc019:~/soft_experiment/10/work$ cd myproject
    c612501080@clpc019:~/soft_experiment/10/work/myproject$ ls
    readme.txt
    c612501080@clpc019:~/soft_experiment/10/work/myproject$ cat readme.txt
    This is readme file.
    c612501080@clpc019:~/soft_experiment/10/work/myproject$ echo Nice to meet you. >> readme.txt
    c612501080@clpc019:~/soft_experiment/10/work/myproject$ cat readme.txt
    This is readme file.
    Nice to meet you.
    c612501080@clpc019:~/soft_experiment/10/work/myproject$ svn diff
    Index: readme.txt
    ===================================================================
    --- readme.txt  (リビジョン 1)
    +++ readme.txt  (作業コピー)
    @@ -1 +1,2 @@
    This is readme file.
    +Nice to meet you.
    c612501080@clpc019:~/soft_experiment/10/work/myproject$ svn commit -m "Add a greeting."
    送信しています              readme.txt
    ファイルのデータを送信しています .
    リビジョン 2 をコミットしました。
    c612501080@clpc019:~/soft_experiment/10/work/myproject$ svn status
    c612501080@clpc019:~/soft_experiment/10/work/myproject$ svn update
    リビジョン 2 です。
    c612501080@clpc019:~/soft_experiment/10/work/myproject$ echo Hello. > hello.txt
    c612501080@clpc019:~/soft_experiment/10/work/myproject$ type hello.txt
    Hello.
    c612501080@clpc019:~/soft_experiment/10/work/myproject$ svn add hello.txt
    A         hello.txt           
    c612501080@clpc019:~/soft_experiment/10/work/myproject$ svn status
    A         hello.txt       
    c612501080@clpc019:~/soft_experiment/10/work/myproject$ svn commit -m "Add hello.txt."
    追加しています              hello.txt
    ファイルのデータを送信しています .
    リビジョン 3 をコミットしました。
    c612501080@clpc019:~/soft_experiment/10/work/myproject$ svn status   
    c612501080@clpc019:~/soft_experiment/10/work/myproject$ svn mkdir src
    A         src
    c612501080@clpc019:~/soft_experiment/10/work/myproject$ cd src
    (テキストエディタでGood.javaを編集)
    c612501080@clpc019:~/soft_experiment/10/work/myproject/src$ cat Good.java 
    public class Good {
        public static void main(String[] args) {
            System.out.println("Good!");
        }
    }
    c612501080@clpc019:~/soft_experiment/10/work/myproject/src$ svn diff
    c612501080@clpc019:~/soft_experiment/10/work/myproject/src$ svn status
    A       .
    ?       Good.java
    c612501080@clpc019:~/soft_experiment/10/work/myproject/src$ svn commit -m "Add an exclamation mark."
    追加しています              .
    Committing transaction...
    リビジョン 4 をコミットしました。
    c612501080@clpc019:~/soft_experiment/10/work/myproject/src$ svn status
    ?       Good.java
    c612501080@clpc019:~/soft_experiment/10/work/myproject/src$ svn add Good.java 
    A         Good.java
    c612501080@clpc019:~/soft_experiment/10/work/myproject/src$ svn commit -m "Add an exclamation mark."
    追加しています              Good.java
    ファイルのデータを送信しています .done
    Committing transaction...
    リビジョン 5 をコミットしました。
    c612501080@clpc019:~/soft_experiment/10/work/myproject/src$ svn move Good.java Main.java
    A         Main.java
    D         Good.java
    c612501080@clpc019:~/soft_experiment/10/work/myproject/src$ ls
    Main.java
    c612501080@clpc019:~/soft_experiment/10/work/myproject/src$ svn status
    D       Good.java
            > moved to Main.java
    A  +    Main.java
            > moved from Good.java
    c612501080@clpc019:~/soft_experiment/10/work/myproject/src$ svn commit -m "Rename it."
    削除しています              Good.java
    追加しています              Main.java
    Committing transaction...
    リビジョン 6 をコミットしました。
\end{verbatim}

\section{感想}
%
IDEは前期の個人開発ではあまり使うことはないかもしれないが, 後期のグループ開発
になると, 必須になるので, 今のうちになれておいて, グループ開発をする時にスムーズ
に使えるようにしておきたい. 授業だけでなく, 社会に出てからもグループで開発すること
が多くなるため, 正しい知識を身に着けておきたい.\\

\begin{thebibliography}{99}
%
% \bibitem{chatGPT} Open AI : ChatGPT \url{https://chatgpt.com/ja-JP} (2026年4月30日閲覧).
% %
% \bibitem{claude} Anthropic : Claude \url{https://claude.ai/new} (2026年5月5日閲覧).
% %
%\bibitem{gemini} Google : Gemini \url{https://gemini.google.com/?hl=ja}
%%
\bibitem{hyuki} 結城 浩 : Subversionの基礎練習 \url{https://www.hyuki.com/techinfo/svninit.html}
%
\end{thebibliography}
\end{document}